\section{Commercial difference-in-differences methods for staggered adoption}
\label{exec18:gp-panel}

A staggered branch rollout poses a different question from the cross-sectional
graphical query. Branches can adopt a policy at different dates, the policy's
effect can change with exposure, and calendar shocks can move all outcomes.
A commercial panel command should therefore be compared on the treatment
history, cohort-time target, admissible controls and uncertainty procedure.
A numerical coefficient without these meanings is insufficient for a banking
decision. The inspected Stata manuals distinguish scalar linear DID and
heterogeneous cohort-time effects; the EViews introduction supplies a more
limited panel-DID capability statement.
\citep{exec18_gp_stata_did,exec18_gp_stata_hdid,exec18_gp_stata_xthdid,exec18_gp_eviews}

The scalar Stata pair is \texttt{didregress} for repeated cross-sections and
\texttt{xtdidregress} for panels. The first fits time and group fixed effects;
the second fits time and panel fixed effects. Their heterogeneous counterparts
are \texttt{hdidregress} and \texttt{xthdidregress}, respectively. The latter
manuals describe regression adjustment, inverse-probability weighting,
augmented weighting and an extended two-way fixed-effects approach. Listing
these alternatives does not make their assumptions or output interchangeable.
The selected heterogeneous panel formulas assume that treatment remains on
after adoption. A history \((0,1,0)\) therefore requires refusal under that
profile or a separately justified switching-treatment design.
\citep{exec18_gp_stata_did,exec18_gp_stata_hdid,exec18_gp_stata_xthdid}

Let \(G_i\) denote branch \(i\)'s first adoption date, with infinity for a
never-treated branch. Write \(Y_{it}(g)\) for its outcome if adoption begins
at \(g\), and \(Y_{it}(\infty)\) for its untreated outcome. A cohort-time
target and one possible aggregation are
\begin{equation}
\label{exec18:gp-cohort-target}
 \theta(g,t)=\mathbb E\{Y_{it}(g)-Y_{it}(\infty)\mid G_i=g\},
 \qquad \Theta_{\omega}=\sum_{g,t\geq g}\omega_{gt}\theta(g,t),
 \quad \omega_{gt}\geq0,\quad\sum_{g,t\geq g}\omega_{gt}=1.
\end{equation}
This notation is our explicit target convention. Its identification needs
the selected comparison assumptions, including consistency, an appropriate
no-anticipation condition, overlap and untreated parallel trends for the
declared comparison population. The inspected manuals allow never-treated
or not-yet-treated controls and document base-time choices for preperiod
effects. Those choices alter which contrasts and support are available.
Not-yet-treated units cease to be eligible untreated comparisons after
adoption. The request must therefore retain the control rule, anticipation
window, base period and the branch/time population before aggregation.
\citep{exec18_gp_stata_hdid,exec18_gp_stata_xthdid}

\subsection{An exact positive-effect rollout can have a negative projection}
\label{exec18:gp-panel-example}

Consider five equally weighted hypothetical branches, \(E_1,E_2,L_1,L_2,N\),
observed at dates \(0,1,2\). Number them \(i=0,\ldots,4\) in that order.
Their untreated outcomes, in thousands of stipulated USD per branch-period,
are \(Y_{it}(\infty)=10i+2t\). The early cohort adopts at date 1, the late
cohort at date 2, and \(N\) never adopts. There is no anticipation, spillover,
attrition or missing period. Treatment adds 1 at the early cohort's first
treated date, 100 at its second, and 1 at the late cohort's treated date.
These are deliberately extreme finite values that make a target error
visible, rather than forecasts of a bank rollout.
\citep{exec18_gp_panel_reference}

\begin{center}
\begin{tabular}{lrrrr}
\hline
Cohort and date & Branches & Effect & ATT weight & TWFE weight\\
\hline
Early, date 1 & 2 & 1 & \(1/3\) & \(5/8\)\\
Early, date 2 & 2 & 100 & \(1/3\) & \(-1/8\)\\
Late, date 2 & 2 & 1 & \(1/3\) & \(1/2\)\\
\hline
\end{tabular}
\end{center}
Here ATT means the average over the six treated branch-periods, so it is
\((2+200+2)/6=34\) thousand USD. Never-treated DID contrasts recover the
three effects exactly: compare each early date with date 0 and the late
treated date with date 1, subtracting the corresponding change in \(N\).
These are finite arithmetic equalities under the stipulated untreated law.
They carry no sampling confidence statement.
\citep{exec18_gp_panel_reference}

Now fit the scalar least-squares projection
\(Y_{it}=\alpha_i+\lambda_t+\beta D_{it}+e_{it}\) to the same 15 rows,
where \(D_{it}=1\{t\geq G_i\}\). In this balanced panel, let bars denote
branch, date and grand means. Removing branch and date fixed effects gives
\begin{equation}
\label{exec18:gp-fwl}
 \widetilde D_{it}=D_{it}-\overline D_i-\overline D_t+\overline D,
 \qquad
 \widehat\beta=\frac{\sum_{i,t}\widetilde D_{it}Y_{it}}
                         {\sum_{i,t}\widetilde D_{it}^{\,2}},
 \qquad \sum_{i,t}\widetilde D_{it}^{\,2}=\frac{16}{15}.
\end{equation}
The treatment functional weights the three cohort-time effects as
\begin{equation}
\label{exec18:gp-negative-weight}
 \widehat\beta=\frac58(1)-\frac18(100)+\frac12(1)
              =-\frac{91}{8}=-11.375,
 \qquad \operatorname{ATT}=34.
\end{equation}
The weights sum to one, but one is negative. The negative projection and
positive treated-period average answer different questions. Exact rational
residualization and a separate solution of the full normal equations agree;
the treatment row of the inverse Gram matrix independently reproduces these
weights. No Stata or EViews engine was used. This example proves a discrepancy
for its specified panel, rather than a universal failure of fixed effects
or a claim about a vendor's undocumented default aggregation.
\citep{exec18_gp_panel_reference}

If the cost is stipulated to be USD 10,000 per treated period, the true
retrospective mean net increment under this law is USD 24,000. Mistaking
the projection for that ATT would instead produce minus USD 21,375.
The latter is an explicitly invalid value calculation. Neither quantity
evaluates a new rollout policy, estimates bank prices or measures ROI.
An investment comparison still needs the applicable population, treatment
costs, feasible alternative and uncertainty about its transported outcomes.
\citep{exec18_gp_panel_reference}

\subsection{Aggregation, trends and uncertainty belong to the request}

The heterogeneous postestimation manual distinguishes aggregation by cohort,
calendar time, exposure time, and cohort/time, and distinguishes
\texttt{weights(timecohort)} from \texttt{weights(cohort)}. The former builds
weights using the subsets used for the individual ATETs; the latter uses
cohort sizes. A branch closure or a changing observed sample can therefore
matter to the interpretation of an aggregate. The adapter should return
the supported cohort-time cells and their denominators first, then the
chosen normalized weights and the resulting estimand. Replacing those
weights by a convenient average can change the banking population.
\citep{exec18_gp_stata_post}

Preperiod diagnostics are useful checks, but cannot validate the unobserved
postperiod counterfactual. For an exact counterexample, let two groups have
identical paths at dates \(-2,-1\). At adoption, suppose the treated group's
true treatment effect is 2, and an otherwise unobserved group-specific
untreated-outcome shock is 3. The observed postperiod DID is 5. A world
with effect 5 and no differential shock has the same observed contrast.
Thus even exact preperiod equality cannot distinguish those worlds without
an additional restriction on the untreated postperiod path. A nonrejected
pretrend test is weaker still. Sensitivity to differential trends and
anticipation must remain part of the business interpretation.
\citep{exec18_gp_panel_reference}

Inference is likewise command-specific. The inspected \texttt{xthdidregress}
manual equates \texttt{vce(robust)} with clustering at the declared treatment
group. That is not a generic promise of independent branch-level rows or
a universal interpretation of every command's robust option. The
postestimation manual disallows its multiplier simultaneous intervals after
TWFE and disallows \texttt{estat aggregation} after bootstrap or jackknife
estimation. A requested interval may target one cell, one fixed aggregate
or simultaneous coverage of many cells; those events differ. A wrapper
must check the selected estimator, clustering, interval target and option
compatibility before returning uncertainty. The five-branch arithmetic
example supplies no cluster asymptotics, multiplier qualification or
confidence coverage. A future generated multi-cluster trial must declare
its sampling law and coverage target separately.
\citep{exec18_gp_stata_xthdid,exec18_gp_stata_post,exec18_gp_panel_reference}

\subsection{What the commercial evidence supports and what to build first}

The acquired Stata originals contain July/August 2026 creation metadata and
StataNow feature flags. They were acquired on 18 October 2026; neither their
metadata nor a displayed copyright authenticates the precise bytes or
entitlements available at the 7 October cutoff. An edition-specific
comparison needs the release and options actually under consideration.
The EViews page, displaying a March 2025 update, states that built-in DID is
provided for panel workfiles and introduces a TWFE and diagnostics chapter.
That introductory page is insufficient evidence of every referenced
algorithm, complete covariance implementation or licensed product edition.
These source limits qualify the historical comparison; they do not imply
that an uninspected vendor lacks a feature.
\citep{exec18_gp_stata_did,exec18_gp_stata_hdid,exec18_gp_stata_xthdid,exec18_gp_stata_post,exec18_gp_eviews}

Our bounded candidate search also retains SAS, IBM SPSS, MATLAB and
commercial frontends to open-source tools as unresolved families.
DID-specific primary originals and exact editions are missing for those
comparisons. Generic panel regression is not sufficient evidence of a
causal DID workflow. The immediate comparison is consequently a documented
Stata baseline and a limited EViews lead, with explicit omissions rather
than a claim of a saturated commercial market. Embedding rights, prices,
support and executable compatibility remain separate selection questions.

The saved local dispatcher probe returned an unsupported disposition for
its dependent-panel DID request. This is an observed limitation of that
request and revision, not evidence that every panel research module is absent.
Chapter~\ref{gap29:public} separately exercises fitted-summary inference through
library, CLI and MCP; it does not qualify a raw-panel DID estimator or erase
the probe's refusal. A proposed extension should admit a
fully declared cohort-time target, preserve the treatment and information
history, report cells and weights, and reject unsupported reversals or
inference options. A licensed vendor comparison can then reproduce the
15-row reference, add unbalanced support and alternative controls, and
examine an independently designed uncertainty trial. Completing this
bounded path would provide a meaningful prototype comparison; it would
not by itself establish general trend validity or bank qualification.
\citep{exec18_gp_dispatch,exec18_gp_panel_reference}
